\subsection{Original Cartier ideals under nonflat normalization}
\begin{lemma}[The original regular principal ideal tensor map]
\label{lem:nonflat-cartier-principal-tensor}
\lean{EnvelopeProperty.IdealTensorBaseChange.toIdealMap_injective_of_regularGenerator,EnvelopeProperty.IdealTensorBaseChange.equivOfRegularGenerator,EnvelopeProperty.IdealTensorBaseChange.equivOfRegularGenerator_tmul}
\leanok
\uses{lem:regular-generator-module-equiv,lem:ideal-tensor-base-change}
Let $R,S$ be arbitrary commutative rings, with an arbitrary original
$R$-algebra structure on $S$. Let $I\subset R$ be the original ideal,
$I=(g)$, where $g$ is a nonzerodivisor in $R$ and its actual image
$g_S$ is a nonzerodivisor in $S$. The existing canonical multiplication
map $S\otimes_R I\longrightarrow IS$, $s\otimes x\mapsto s x$,
is injective and hence a genuine $S$-linear isomorphism. No flatness,
Noetherian, reducedness, finiteness or field hypothesis is imposed.
\end{lemma}
\begin{proof}
Use the existing multiplication-by-$g$ isomorphism $R\simeq_R I$,
base change it to $S$, and compose with $S\otimes_R R\simeq_S S$.
The canonical tensor map in these coordinates is multiplication by the
actual $g_S$. Indeed on every pure tensor $s\otimes r$ both composites
have included value $s\,\operatorname{im}(r)\,g_S$; tensor induction proves
whole-map equality. Multiplication by $g_S$ is injective by regularity,
so the canonical map is injective. Its existing unconditional surjectivity
onto the actual extended ideal gives the stated isomorphism. All rings,
algebra actions, ideals and canonical maps are unchanged; nilpotents
outside the regular equation are retained.
\end{proof}

\begin{lemma}[Dominant Cartier pullback without flatness]
\label{lem:nonflat-dominant-cartier-comap}
\lean{AlgebraicGeometry.Scheme.IdealSheafData.IsInvertible.comap_of_dominant_integral}
\leanok
\uses{scheme:genuine-ideal-pullback-stalk}
Let $f:X\to Y$ be an actual dominant morphism with $X$
integral and $Y$ reduced. For every actual effective Cartier ideal
$J$ on $Y$, its original geometric ideal comap $J\mathcal O_X$ is
an actual effective Cartier ideal. No quasi-compactness, flatness, finite type,
Noetherian, separation, projectivity or normality is assumed.
\end{lemma}
\begin{proof}
Fix an original $x\in X$. Choose an affine Cartier equation $g$ for $J$
on an original affine neighborhood $U$ of $f(x)$; $g$ is regular there.
Choose an original affine neighborhood $V$ of $x$ inside $f^{-1}U$.
Since $X$ is integral, its nonempty open $V$ is dense; thus $V\to X$
and its composite $V\to Y$ are dominant. Factoring through the original
open immersion $U\to Y$ shows that the actual affine restriction
$f_{UV}:V\to U$ is dominant. Reducedness of $U$ implies scheme-theoretic
dominance by the existing Mathlib theorem. The morphism $f_{UV}$ between
affines is automatically quasi-compact, so its actual section map
$\Gamma(U)\to\Gamma(V)$ is injective. The original $U$ is nonempty
because it contains $f(x)$, so its regular equation is nonzero.
Hence the actual $g_V$ is nonzero;
it is regular in the original domain $\Gamma(V)$. The existing geometric
ideal-comap formula gives $(J\mathcal O_X)(V)=(g_V)$. These original
neighborhoods and equations prove the original effective Cartier definition.
\end{proof}

\begin{corollary}[The actual normalization retains original Cartier ideals]
\label{lem:nonflat-normalization-cartier-comap}
\lean{EnvelopeProperty.IntegralSchemeNormalization.cartier_comap}
\leanok
\uses{lem:nonflat-dominant-cartier-comap,scheme:actual-integral-normalization-surjective}
For any integral original scheme $Y$, let $n:Y^\nu\to Y$ be the EXISTING
relative normalization in its original function field, with its EXISTING
projection. Every original effective Cartier ideal $J$ has effective Cartier
comap along this actual $n$. No finite normalization, finite type,
Noetherian, excellent, field, quasi-compact or separation assumption on $Y$
is imposed.
\end{corollary}
\begin{proof}
The existing actual generic-field spectrum is integral, so Mathlib's
relative normalization is integral. The original projection has already proved surjectivity for this exact normalization; thus it is dominant. The original target is reduced because
it is integral. Apply the preceding genuine theorem to $n$.
This proves Cartier preservation, not a normalized polarized-model,
relative-nef approximation or the LocalPositivity theorem. The canonical
nonflat module-sheaf comparison is a subsequent separate obligation.
\end{proof}
